\scanpage{127}
% line head | text | 76
\mspage{76}
% line 1–3 | diagram | (*) intersection, U, U prime, X
\diagram{\begin{tikzcd}[ampersand replacement=\&,row sep=small,column sep=small]
\& \& U\cap U'\arrow[dl,hook',"i'_0"']\arrow[dr,hook,"i_0"] \&\\
{(\unsure{*|\alpha})} \& U\arrow[dr,hook,"i"'] \& \& U'\arrow[dl,hook',"i'"]\\
\& \& X \&
\end{tikzcd}}
% line 4 | text | Alors
Alors
% line 5 | display | i_0 in W
\[
i_0\in W\Longleftrightarrow i\in W
\]
% line 6 | text | et symétriquement
(et symétriquement
% line 7 | display | i_0 prime in W
\[
i'_0\in W\Longleftrightarrow i'\in W\ \bigr).
\]
% line 8–9 | text | NB Par dualité
\uline{NB} Par dualité cet axiome est équivalent à l'axiome similaire
avec des \uline{fermés} $T,T'$ de $X$ de réunion $X$.

% line 10–11 | text | Corollaire
\uline{Corollaire}\quad \inserted{\fbox{Supposons $W(7\mathrm{bis})$.}} Considérons les conditions \unsure{suiv.} \deleted{\unsure{de}} sur
le carré précédent

% line 12 | text | a) U intersection U prime
a) $U\cap U'$ \inserted{$W$-}asphérique

% line 13 | text | b) U, U prime
b) $U,U'$ $W$-asphériques

% line 14 | text | c) X
c) \unsure{$X$} $W$-asphérique

% line 15 | text | Alors (a et b)
Alors (a et b) \unsure{impliquent} c, et (b et c) implique a.

% line 16–22 | text | NB Moralement
\unsure{\{} \uline{NB} Moralement, il faudrait \inserted{aussi} que (c \deleted{$+$} et a) implique
\deleted{\unsure{aussi}} b, i.e. si $U\cap U'$ et $X$ sont $W$-asphériques,
il en est de même de $U$ et de $U'$. Mais
il ne semble pas que l'axiome du carré cocartésien
permette de le prouver. Par contre, l'axiome
qui suit impliquera le corollaire précédent sous sa
forme complète.

% line 23–24 | text | W, Axiome du cube, Lemme des cinq
\deleted{\uline{\unsure{NB}}} $W(\unsure{7|8})$ (\inserted{``Axiome du cube'', ou} ``Lemme des cinq'') Soit un diagramme
commutatif cubique
% line 25–28 left | diagram | cancelled drawing, intersection at top through X1, lower part begun
\diagram{\deleted{\begin{tikzcd}[ampersand replacement=\&,row sep=tiny,column sep=tiny,execute at end picture={
\draw[thin] (current bounding box.north) -- (current bounding box.west) -- (current bounding box.north east) -- (current bounding box.south west) -- (current bounding box.east) -- (current bounding box.south);
}]
\& \& U_1\cap U'_1\arrow[dl,hook',"i'_0"']\arrow[dr,hook,"i_0"] \&\\
\& U_1\arrow[dr,"i"']\arrow[ddl] \& \& U'_1\arrow[dl,"i'"]\\
\& U_2\cap U'_2\arrow[dl] \& X_1 \&\\
U_2 \& \& \&
\end{tikzcd}}}
% line 25–28 right | diagram | live cube comparing intersections, opens and unions
\diagram{\begin{tikzcd}[ampersand replacement=\&,row sep=small,column sep=small]
U_1\cap U'_1\arrow[rr]\arrow[dr]\arrow[dd,"f_{U_1\cap U'_1}"'] \& \& U_1\arrow[dr]\arrow[dd,"f_{U_1}" description] \&\\
\& U'_1\arrow[rr]\arrow[dd,"f_{U'_1}"] \& \& X_1\arrow[dd,"f"]\\
U_2\cap U'_2\arrow[rr]\arrow[dr] \& \& U_2\arrow[dr] \&\\
\& U'_2\arrow[rr] \& \& X_2
\end{tikzcd}}
% line 29 | text | représentant un morphisme
représentant un morphisme (flèches verticales) d'un
% archive sheet number: 127
